\documentclass[11pt, a4paper]{article}
\usepackage{fontspec}
\usepackage[turkish]{babel}
\setmainfont{Latin Modern Roman}
\usepackage[margin=2.5cm]{geometry}
\usepackage{hyperref}
\usepackage[shortlabels]{enumitem}
\usepackage{xcolor}
\usepackage{longtable}
\usepackage{booktabs}
\usepackage{tabularx}
\usepackage{listings}

\lstset{
    basicstyle=\ttfamily\small,
    frame=single,
    breaklines=true,
    columns=fullflexible,
}

\hypersetup{
    colorlinks=true,
    linkcolor=blue,
    urlcolor=cyan,
    pdftitle={AC2 Interface Requirements Specification},
}

\title{\Huge \textbf{Interface Requirements Specification}\\
       \LARGE \textbf{AC2 Telemetry Protocol}}
\author{\Large \textbf{Sürüm:} 2.0.0 \hfill \textbf{Doc ID:} AC2-IRS-001}
\date{\today}

\begin{document}
\maketitle
\tableofcontents
\newpage

\section{Introduction}
\subsection{Purpose}
Bu IRS, STM32 Edge Node ile Gateway Bridge arasındaki seri haberleşme arayüzünü (UART) ve Gateway ile Web UI arasındaki WebSocket arayüzünü byte-level netlikte tanımlar. Doküman, AC2-SRD-001 [R4]'ün CPR bölümünü yerine alır ve genişletir.

\subsection{Interface Overview}
\begin{itemize}
    \item \textbf{I/F-A:} STM32 Edge Node <--> Gateway (UART, 115200 8N1, AC2 framed binary).
    \item \textbf{I/F-B:} Gateway <--> Web UI (WebSocket Secure, JSON messages).
\end{itemize}

\section{Physical Layer (I/F-A)}

\begin{tabularx}{\textwidth}{|l|X|}
\hline
\textbf{Parameter} & \textbf{Value} \\
\hline
Signal Levels & 3.3V CMOS (STM32) / TTL (USB-UART adapter) \\
Baud Rate & 115200 \\
Data Bits & 8 \\
Parity & None \\
Stop Bits & 1 \\
Flow Control & None (RTS/CTS kapalı) \\
Endianness & Little-endian (wire) \\
Max Frame Size & 64 byte \\
Idle Detection & UART IDLE IRQ (DMA boundary) \\
\hline
\end{tabularx}

\section{AC2 Frame Format (I/F-A)}

\subsection{Frame Structure}
\begin{lstlisting}
+---------+---------+----------+----------+---------+---------+----------+---------+
| SYNC    | VERSION | SEQ_NUM  | LENGTH   | CMD_ID  | PAYLOAD | HMAC-8   | CRC-16  |
| 1 byte  | 1 byte  | 4 bytes  | 1 byte   | 1 byte  | 0-48 B  | 8 bytes  | 2 bytes |
+---------+---------+----------+----------+---------+---------+----------+---------+
  0xAA      0x02      LE u32     N         (Table)   variable  (trunc)    CRC-16-CCITT
\end{lstlisting}

\subsection{Field Semantics}
\begin{tabularx}{\textwidth}{|l|l|X|}
\hline
\textbf{Field} & \textbf{Size} & \textbf{Description} \\
\hline
SYNC    & 1 B & Sabit \texttt{0xAA}. Senkronizasyon başlangıcı. \\
VERSION & 1 B & Protocol sürümü. v2.0.0 için \texttt{0x02}. \\
SEQ\_NUM & 4 B LE & Monotonic artan frame numarası. Replay koruması (SR-02). \\
LENGTH  & 1 B & PAYLOAD uzunluğu (0..48). HMAC ve CRC bu sayıya dahil değildir. \\
CMD\_ID & 1 B & Komut/olay ID (bkz. 4.). \\
PAYLOAD & N B & Komuta özgü binary data. Endian: LE. \\
HMAC-8  & 8 B & HMAC-SHA-256'nın ilk 8 byte'ı. \texttt{CMD\_ID + PAYLOAD} üzerinden; komut yönünde zorunlu, telemetry yönünde tüm sıfır (ignore). \\
CRC-16  & 2 B LE & CRC-16/CCITT-FALSE (poly 0x1021, init 0xFFFF, no reflect, xorout 0x0000). SYNC dahil tüm byte'lar üzerinden. \\
\hline
\end{tabularx}

\subsection{Byte Stuffing / Escaping}
Payload içinde \texttt{0xAA} byte'ı geçebilir (binary data). Parser'ın SYNC arayarak resync olmasını engellemek için \textbf{length-prefixed framing} yaklaşımı kullanılır: SYNC'i gören alıcı, hiç escaping yapmaksızın \texttt{VERSION+SEQ+LENGTH} header'ını okur ve frame boyunu deterministik hesaplar. Eğer CRC fail ise tek byte ilerleyerek yeniden SYNC aramaya başlar.

\subsection{Maximum Frame Computation}
\texttt{1 + 1 + 4 + 1 + 1 + 48 + 8 + 2 = 66 byte}. Pratik limit 64 byte payload $\leq$ 48 B ile korunur. Daha büyük veri için fragmentation protokolü (v3.0 roadmap) gerekir.

\section{Command Registry}

\begin{longtable}{|l|l|l|l|}
\hline
\textbf{CMD\_ID} & \textbf{Name} & \textbf{Direction} & \textbf{Payload} \\
\hline
\endhead
0x01 & CMD\_GET\_VERSION      & Host -> Edge & — \\
0x02 & CMD\_VERSION\_REPORT   & Edge -> Host & \texttt{u8 maj, u8 min, u8 patch, u8[8] git\_sha, u32 build\_ts} \\
0x10 & RESERVED\_LEGACY\_SET\_STATE & — & Deprecated; UI/Gateway artık state set etmez \\
0x11 & CMD\_MANUAL\_LOCK      & Host -> Edge & — (state SEARCH/TRACK'tayken idempotent) \\
0x12 & CMD\_RESET\_ACK        & Host -> Edge & — (FAIL-SAFE'ten çıkış için) \\
0x13 & CMD\_SYSTEM\_RESET     & Host -> Edge & — (ACK sonrası gecikmeli MCU reboot) \\
0x20 & CMD\_REPORT\_STATE     & Edge -> Host & \texttt{u8 state, u8 prev\_state, u32 uptime\_ms} \\
0x21 & CMD\_TELEMETRY\_TICK   & Edge -> Host & \texttt{u8 state, u16 cpu\_load\_x10, u16 free\_stack\_min\_words, u8 hb\_miss\_count} \\
0x22 & CMD\_TASK\_LIST        & Edge -> Host & \texttt{u8 count, packed task entries[14]} \\
0x23 & CMD\_RANGE\_SCAN\_REPORT & Edge -> Host & \texttt{u8 angle, u16 distance\_cm, u8 flags, u16 threshold\_cm} \\
0x24 & CMD\_BOOT\_REPORT      & Edge -> Host & \texttt{u32 reset\_reason\_bits} \\
0x30 & CMD\_FAULT\_REPORT     & Edge -> Host & \texttt{u8 fault\_code, u16 ctx, u32 reset\_reason\_bits} \\
0x31 & CMD\_AUDIT\_EVENT      & Edge -> Host & \texttt{u8 event\_code, u16 count} \\
0x50 & CMD\_CREATE\_TASK      & Host -> Edge & \texttt{u8 task\_type, u8 param} \\
0x51 & CMD\_DELETE\_TASK      & Host -> Edge & \texttt{u8 slot\_index} \\
0x60 & CMD\_DETECTION\_RESULT & Host -> Edge & \texttt{u8 class\_id, u8 confidence\_pct} \\
0x80 & CMD\_ACK               & Edge -> Host & \texttt{u32 echoed\_seq} \\
0x81 & CMD\_NACK              & Edge -> Host & \texttt{u32 echoed\_seq, u8 err\_code} \\
0x99 & CMD\_HEARTBEAT         & Bidirectional & \texttt{u32 uptime\_ms} (TX), \texttt{u32 host\_ts} (RX) \\
\hline
\end{longtable}

\subsection{Audit Event Codes}
\begin{tabularx}{\textwidth}{|l|l|X|}
\hline
\textbf{Code} & \textbf{Name} & \textbf{Meaning} \\
\hline
0x01 & AUDIT\_BOOT        & Edge booted; \texttt{count} reserved. \\
0x02 & AUDIT\_TASK\_CREATE & User task created; \texttt{count}=slot index. \\
0x03 & AUDIT\_TASK\_DELETE & User task deleted; \texttt{count}=slot index. \\
0x04 & AUDIT\_RANGE\_LOCK  & Range scan locked; \texttt{count}=angle degrees. \\
0x05 & AUDIT\_VISION\_HIT  & Vision detection crossed threshold; \texttt{count}=confidence percent. \\
0x06 & AUDIT\_SYSTEM\_RESET & System reset command accepted; \texttt{count} reserved. \\
0x07 & AUDIT\_FAIL\_SAFE   & Fail-safe entered; \texttt{count} reserved. \\
\hline
\end{tabularx}

\subsection{Error Codes (CMD\_NACK.err\_code)}
\begin{tabularx}{\textwidth}{|l|l|X|}
\hline
\textbf{Code} & \textbf{Name} & \textbf{Description} \\
\hline
0x01 & ERR\_INVALID\_CMD        & Bilinmeyen CMD\_ID. \\
0x02 & ERR\_INVALID\_PAYLOAD    & Payload uzunluğu veya içeriği geçersiz. \\
0x03 & ERR\_INVALID\_TRANSITION & State transition tablosunda yok. \\
0x04 & ERR\_AUTH\_FAIL          & HMAC doğrulaması başarısız. \\
0x05 & ERR\_REPLAY              & SEQ\_NUM $\leq$ son görülen. \\
0x06 & ERR\_RATE\_LIMITED       & Token bucket dolu. \\
0x07 & ERR\_FAIL\_SAFE\_LOCK    & Sistem FAIL-SAFE modunda; komut reddedildi. \\
0x08 & ERR\_BUSY                & Geçici doluluk; sonra tekrar dene. \\
\hline
\end{tabularx}

\section{Timing Requirements}

\begin{tabularx}{\textwidth}{|l|l|X|}
\hline
\textbf{Parameter} & \textbf{Value} & \textbf{Notes} \\
\hline
Heartbeat period (Edge $\to$ Host) & 1000 ms $\pm$ 50 ms & SAF-04 \\
Heartbeat timeout (Host$\to$Edge loss) & 3 s & Degraded trigger \\
FAIL-SAFE trigger & 5 s & SAF-04 \\
Telemetry tick period & 100 ms & Throughput baseline \\
Max ACK latency (Edge response to command) & 100 ms & \\
\hline
\end{tabularx}

\section{WebSocket Interface (I/F-B)}

\subsection{Endpoint}
\begin{itemize}
    \item URL: \texttt{wss://<host>:8443/ws}
    \item Subprotocol: \texttt{ac2.v2}
    \item Auth: \texttt{Sec-WebSocket-Protocol: bearer.<JWT>, ac2.v2}
\end{itemize}

\subsection{Message Envelope}
Tüm mesajlar JSON'dır ve şu envelope'u taşır:
\begin{lstlisting}
{
  "v":    2,
  "type": "<event|command|ack|error>",
  "ts":   "2026-04-21T18:57:23.412Z",
  "seq":  12345,
  "data": { /* type-specific */ }
}
\end{lstlisting}

\subsection{Event Types (Gateway -> UI)}
\begin{tabularx}{\textwidth}{|l|X|}
\hline
\textbf{type} & \textbf{data schema} \\
\hline
\texttt{state}      & \{ state, prev\_state, uptime\_ms \} \\
\texttt{telemetry}  & \{ state, cpu\_load, free\_stack, hb\_miss \} \\
\texttt{fault}      & \{ code, ctx, reset\_reason, severity \} \\
\texttt{audit}      & \{ event\_code, count \} \\
\texttt{connection} & \{ status: "connected"|"degraded"|"disconnected" \} \\
\hline
\end{tabularx}

\subsection{Command Types (UI -> Gateway)}
\begin{tabularx}{\textwidth}{|l|X|}
\hline
\textbf{type} & \textbf{data schema} \\
\hline
\texttt{set\_state}   & \{ target: "IDLE"|"SEARCH"|"TRACK" \} \\
\texttt{manual\_lock} & \{\} \\
\texttt{reset\_ack}   & \{\} \\
\texttt{get\_version} & \{\} \\
\hline
\end{tabularx}
Her komut Gateway tarafından bir \texttt{ack} veya \texttt{error} ile cevaplanır (echoed \texttt{seq}).

\subsection{Validation}
Gateway tüm incoming mesajları Zod schema'larına göre doğrular. Şema dışı mesaj drop edilir ve \texttt{audit} event'i yayılır (SR-08).

\section{Interface Versioning}
\begin{itemize}
    \item \texttt{VERSION} byte'ı MAJOR değişiklikte artar. MINOR eklemeler geriye dönük uyumludur.
    \item Handshake: Gateway bağlandığında \texttt{CMD\_GET\_VERSION} ile edge firmware sürümünü öğrenir, uyumsuzluğu UI'ya \texttt{fault} event'i olarak bildirir.
\end{itemize}

\section{Test Vectors}
\subsection{Example: CMD\_SET\_STATE(SEARCH)}
\begin{lstlisting}
Direction: Host -> Edge
Seq: 0x00000042
PSK:  (32-byte demo key, see tests/fixtures)

Plaintext body (cmd + payload):
  10 01

HMAC-SHA-256 over body (trunc 8):
  A7 3F 91 22 4B 88 0E 1C

Frame bytes:
  AA 02 42 00 00 00 02 10 01 A7 3F 91 22 4B 88 0E 1C <CRC_LO> <CRC_HI>

CRC-16-CCITT-FALSE over SYNC..HMAC = 0x5F A3 (example)
Final bytes:
  AA 02 42 00 00 00 02 10 01 A7 3F 91 22 4B 88 0E 1C A3 5F
\end{lstlisting}

\subsection{Example: CMD\_REPORT\_STATE (SEARCH)}
\begin{lstlisting}
Direction: Edge -> Host
Seq: 0x00000101
HMAC: all zeros (telemetry direction)

Frame:
  AA 02 01 01 00 00 06 20 01 00 F0 1A 00 00 00 00 00 00 00 00 00 <CRC_LO> <CRC_HI>
   SYN VER  <-- SEQ -->  LEN CMD <-- payload(6) -->   <-- HMAC(8 zero) -->  CRC16
\end{lstlisting}

\end{document}
